\section{A Wider Covariance Box and Transfer to Averaged TD}
\label{sec:td-extension}
We establish main Theorem~\ref{thm:finite-robust-separation} over $\theta_0\in[0.04,0.06]$, $\theta_1\in[1.8,2.2]$ and $\lambda_j\in[0.15,0.25]$, with prior $\pi\in[0.45,0.55]$.
The reference detector and its eighteen observations remain fixed.

\subsection{Direct error control for the frozen detector}
Let $K_0,K_1$ be its reference contrast covariances and $Q=K_0^{-1}-K_1^{-1}\succ0$.
Their generalized eigenvalues satisfy $\ell_i\ge\ell_*=175/19$ as in Section~\ref{sec:finite-robust}.
The high decision is $y^\top Qy\ge\sum_i\log\ell_i$.
Throughout the enlarged box, the true contrast covariances obey
\[
\widetilde K_0\preceq c_0K_0,\quad \widetilde K_1\succeq c_1K_1,
\quad c_0=9/8,\ c_1=4/5.
\]
For the low model, the common covariance is at most $4/3$ times its reference: the variance ratio is at most $6/5$ and the spectral ratio at most $10/9$.
The reference common covariance is bounded by $(3/40)I$, while private noise contributes $I/8$.
Its fraction of the total covariance is at most $3/8$, giving total factor $1+(1/3)(3/8)=9/8$.
For the high model, the variance and spectral ratios are each at least $9/10$; adding private noise preserves a total lower factor $81/100>4/5$.
Congruence by the contrast matrix preserves both comparisons.

The Gaussian quadratic-form transform and Markov's inequality give
\begin{align*}
p_0&\le\prod_{i=1}^{17}\ell_i^{-1/4}[1-c_0/2+c_0/(2\ell_i)]^{-1/2},\\
p_1&\le\prod_{i=1}^{17}\ell_i^{1/8}[1+c_1(\ell_i-1)/4]^{-1/2}.
\end{align*}
The exponential parameters are $1/4$ for the outside event and $1/8$ for the inside event.
Covariance order bounds the positive-quadratic transform in the first case and the negative-quadratic transform in the second.
The first transform is finite because $1-c_0/2>0$.
Logarithmic differentiation shows both factors decrease for $\ell\ge175/19$.
Their fourth and eighth powers at $\ell_*$ satisfy the exact rational inequalities
\begin{align*}
\ell_*^{-17}[1-c_0/2+c_0/(2\ell_*)]^{-34}&<(3/100)^4,\\
\ell_*^{17}[1+c_1(\ell_*-1)/4]^{-68}&<(3/100)^8.
\end{align*}
Hence $p_0,p_1<0.03$ throughout the enlarged box.
This controls the unchanged detector directly instead of adding a total-variation perturbation to its reference error bound.

\subsection{An explicit averaged-TD boundary bound}
Consider scalar TD(0) with constant feature, discount $\gamma$, $\mu=1-\gamma>0$ and contraction $\beta=1-\eta\mu\in(0,1)$.
For centered Markov reward noise $\bar X_t$, its critic error satisfies $e_t=\beta e_{t-1}+\eta\bar X_t$.
The average of the first $N$ iterates has initialization coefficient and noise weights
\[
d_N=\frac{\beta(1-\beta^N)}{N(1-\beta)},\quad
w_t=\frac{1-\beta^{N-t+1}}{\mu N}.
\]
Let $C_N$ be the observation covariance, with nonnegative entries and row sums at most $C$.
One may take $C=1+\theta(1+\lambda)/(1-\lambda)$ for any resource schedule, since participation and physical gaps are at least one.
Put $z_t=\beta^{N-t+1}$ and $S_N=\mathbf1^\top C_N\mathbf1$.
Positive semidefiniteness gives
\begin{align*}
w^\top C_Nw
&=\frac{S_N-2z^\top C_N\mathbf1+z^\top C_Nz}{\mu^2N^2}\\
&\ge\frac{S_N-2C\beta/(1-\beta)}{\mu^2N^2}.
\end{align*}
Thus the unweighted resource envelope transfers with an explicit $O(N^{-2})$ correction.
The upper bound $w_t\le1/(\mu N)$ gives variance at most $v/(\mu^2N)$ for a constant program with long-run noise coefficient $v$.

Use $\mu=1/2$, $\eta=1/10$, $\beta=19/20$ and zero update delay.
For $N\le512$, monotonicity of $d_N$ bounds worst-initialization risk below by $d_{512}^2$.
For $N>512$, put $b_*=\beta^2(1-\beta^{512})^2/(1-\beta)^2$, so $N^2d_N^2\ge b_*$.
Let $v_-(n)$ be the resource-envelope coefficient using the separate lower persistence of each model.
Set $a_-=2\sum_j\pi_j\theta_j\lambda_{j,-}$ and $C_+=1+\sum_j\pi_j\theta_j(1+\lambda_{j,+})/(1-\lambda_{j,+})$.
Every predetermined TD schedule has risk at least the minimum of $d_{512}^2$ and
\[
\frac{b_*+4\{nv_-(n)-a_--38C_+\}}{n^2},\quad 513\le n\le16B/9.
\]
On each resource segment the numerator is affine with positive slope; its ratio to $n^2$ has no interior minimum.
Evaluating segment endpoints therefore covers all stopping lengths.
For paid constant-program continuation of length $N$, use
\[
u_{j,a}=361/N^2+4v_{j,a}/N,\quad
\kappa_a=4\rho_a^2/[N^2(1-\rho_a)^2].
\]
These bound stationary risk and squared probe-state transfer, respectively.
Evenness of the detector cancels the unknown-initialization cross term.
Insert these expressions into the paid upper bound of Section~\ref{sec:finite-robust}, with directional error $0.03$.
For TD, $N=\min\{\lfloor(16B-288)/(8+q)\rfloor,\lfloor(4B-18)/b\rfloor\}$; for the running mean retain its eight withheld updates.

\subsection{Uniform finite-budget comparisons}
Partition each persistence interval into sixteen closed intervals of width $1/160$.
Use lower endpoints in the open-loop covariance envelope and upper endpoints in the paid risk and TD row-sum bound.
Evaluate both variance endpoints and both prior endpoints in every cell.
For fixed persistence cell, the lower bound is separately concave in each variance and the prior, being a minimum of affine functions.
The paid upper bound is separately convex, since it uses maxima of affine branch risks.
Endpoint inequalities therefore imply the comparison throughout each cell and hence the entire box.
This is an interval proof over a continuum, not a parameter-grid simulation.

\begin{table}[t]
\centering
\caption{Uniform analytical risk-reduction guarantees over the enlarged covariance box, against every predetermined resource schedule.}
\label{tab:wide-td}
\begin{tabular}{rrr}
\toprule
$B$ & Running mean, $D=8$ & Averaged TD, $D=0$\\
\midrule
8192 & $\ge2.54\%$ & $\ge1.86\%$\\
32768 & $\ge2.75\%$ & $\ge2.58\%$\\
\bottomrule
\end{tabular}
\end{table}

At $B=8192$ the respective unrounded guarantees exceed $0.0254335$ and $0.0186558$.
Exact rational cross multiplication proves the rounded main-theorem constants.
The TD result covers a constant-feature affine critic with exponential filtering; general random TD Jacobians require a different weight argument.
These calculations alter no registered empirical comparison or its decisions.
